\section*{Ethical and Privacy Considerations}

LightMem-Ego relies on egocentric visual and audio streams, which may contain sensitive information about users and nearby bystanders, including private conversations, faces, documents, locations, screens, and daily routines. 
Because the system converts these streams into persistent and queryable memory, privacy risks can arise from both raw data and derived representations such as transcripts, event summaries, embeddings, and long-term semantic memories.

The current prototype is intended as a research demonstration and has not yet implemented a complete privacy-preserving pipeline. 
It does not automatically redact sensitive content, manage bystander consent, enforce fine-grained access control, or provide mature retention and deletion policies. 
Future work will integrate privacy protection into the memory lifecycle, including on-device preprocessing, selective capture, sensitive-content filtering, encrypted storage, user-controlled memory editing and deletion, and privacy-aware consolidation that avoids promoting incidental sensitive information into long-term memory.